\section{Verification, provenance, and reproducibility}
\label{sec:verification}

\subsection{Independent mathematical checks}
The report was prepared with AI assistance. The independent reconstructions and reviews in this archive were performed by separate collaborating agents; they are not external journal refereeing.

The torsion proof was independently reconstructed and reviewed, including:
arbitrary \(\Cp\)-roots; repeated roots; the case \(p=2,h_0=1\); integer budget boundaries; off-shell roots; nonattained coefficient infima; arbitrary finite dimension; bounded-series truncation; aggregate packets; and the Lubin--Tate normalization. The complete review is included.

The calculations supplement the proofs. They use exact integer, rational, or quadratic-field arithmetic and do not rely on floating-point estimates.

\begin{center}
\small
\begin{tabularx}{\textwidth}{@{}p{.34\textwidth}X@{}}
\toprule
Check & Scope and observed result\\
\midrule
Aggregate recurrence & 800 greedy shell recurrences and 800 floor identities: PASS.\\
Boundary thresholds & 1,440 rational budget cases, including equality: PASS.\\
Sharp constructions & 55 constructions, 22,473 primitive evaluation points: PASS.\\
Independent shell minimax & All 6,974 marked multisets in seven small shells and 27 budget cases: PASS.\\
\(A_5\) certificate & Exact character arithmetic in \(\mathbf Q(\sqrt5)\), actual subgroup induction, all 59 subgroups, maximal restrictions, and six-ray inequalities: PASS.\\
Source audit & Both editions scanned; file hashes and file-level import closure recorded. No independent kernel replay.\\
Prior-source provenance & All 31 archive lengths and Git blob hashes matched the pinned directory.\\
\bottomrule
\end{tabularx}
\end{center}

\subsection{Archive contents}
\begin{itemize}
\item \path{report.tex}, \path{report.pdf}, and \path{sections/}: complete manuscript and sources.
\item \path{code/verify_aggregate_minimax.py}: shell recurrence, equality, threshold, and extremizer checks.
\item \path{code/verify_shell_minimax_exhaustive.py}: an independent exhaustive maximin computation over small marked multisets.
\item \path{code/verify_a5_tensor.py}: exact character certificate and subgroup enumeration.
\item \path{code/make_figures.py}: exact bound table and the comparison plot.
\item \path{code/scan_source.py}: reproducible source inventory and lexical audit; requires the frozen upstream checkout.
\item \path{data/}: exact verification results and bound data.
\item \path{evidence/}: source audit, prior-corpus comparison, immutable source manifest, and independent proof review.
\item \path{formalization/plan.md}: a proposed Lean decomposition and the precise unformalized boundaries.
\item \path{build.sh}, \path{README.md}, and \path{SHA256SUMS}: reproduction instructions and archive integrity.
\end{itemize}

\subsection{Reproduction}
From the extracted directory, run:
\begin{lstlisting}
python3 code/verify_aggregate_minimax.py
python3 code/verify_shell_minimax_exhaustive.py
python3 code/verify_a5_tensor.py
python3 code/make_figures.py
bash build.sh
\end{lstlisting}
The three proof-check scripts use the Python standard library. Plotting additionally uses Matplotlib. PDF compilation uses a standard LaTeX installation and \texttt{latexmk}.

The scripts' existing exact assertions are the main reproducible checks. JSON files record the observed outputs. A supplied script is not a formal proof in Lean, and the manuscript makes no such claim.

To reproduce the upstream Lean theorem checks separately, obtain the pinned source and follow its README:
\begin{lstlisting}
cd grh
lake exe cache get
lake build
lake env lean scripts/PrintAxioms.lean
\end{lstlisting}
The upstream build was not run here. No GitHub branch, issue, pull request, or repository file was modified.

\subsection{Interpretation of the result}
The sharp torsion formula is a completed mathematical theorem, with an explicit optimality example, and improves the same objective studied in an inspected predecessor. The tensor identities are complete finite-character statements with a rigorous analytic interpretation. The new nonabelian prime-sum estimate remains a research target. These categories should remain separate when incorporating the report into ProveIt.

